\subsection{The Hille-Yosida Theorem}\label{sec:hille-yosida-theorem}

\subsubsection{The Hille-Yosida Theorem}\label{subsec:hille-yosida-theorem}

In the preceding section, we showed that every strongly continuous semigroup has an infinitesimal generator and also how the semigroup can be used to solve initial-value problems. The most important question from the PDE point of view is how to recognize those operators which are infinitesimal generators of a $C_0$-semigroup. The Hille-Yosida theorem provides an answer to that question.

\begin{Theorem}{Hille-Yosida}
\label{thm:hille-yosida-generation}
Let $A$ be an operator in the Banach space $X$. Then $A$ is the infinitesimal generator of a $C_0$-semigroup $T(t)$ satisfying $\|T(t)\|\leq M\exp(\omega t)$ if and only if the following two conditions hold:\begin{enumerate}\item\label{hyp:hille-yosida-dense-closed} $\mathcal{D}(A)$ is dense and $A$ is closed,\item\label{hyp:hille-yosida-resolvent-powers} Every real number $\lambda>\omega$ is in the resolvent set of $A$ and $\|R_\lambda(A)^n\|\leq M/(\lambda-\omega)^n$ for every $n\in\mathbb{N}$.\end{enumerate}
\end{Theorem}

It will be clear from the proof below that actually every complex number with $\operatorname{Re}\lambda>\omega$ is in the resolvent set and that $\|R_\lambda(A)^n\|\leq M/(\operatorname{Re}\lambda-\omega)^n$.

\begingroup
\renewcommand{\endProof}{\par}
\begin{Proof}
We already know that the first condition is necessary. To see the necessity of the second, recall that for any $C_0$-semigroup there are constants $M$ and $\omega$ such that $\|T(t)\|\leq M\exp(\omega t)$. Let us now consider the integral
\begin{equation}\label{eq:hille-yosida-resolvent-power-laplace-integral}
I_n(\lambda)x=\dfrac{1}{(n-1)!}\int_0^\infty t^{n-1}e^{-\lambda t}T(t)x\,\mathrm{d}t.
\end{equation}
Clearly, $I_n(\lambda)$ is well defined for $\operatorname{Re}\lambda>\omega$ and
\begin{equation}\label{eq:hille-yosida-laplace-integral-norm-bound}
\|I_n(\lambda)\|\leq\dfrac{1}{(n-1)!}\int_0^\infty t^{n-1}e^{-\operatorname{Re}\lambda t}Me^{\omega t}\,\mathrm{d}t=\dfrac{M}{(\operatorname{Re}\lambda-\omega)^n}.
\end{equation}
If $x\in\mathcal{D}(A)$, we find, for $n>1$,\[\begin{aligned}
I_n(\lambda)Ax&=\dfrac{1}{(n-1)!}\int_0^\infty t^{n-1}e^{-\lambda t}T(t)Ax\,\mathrm{d}t\\
&=\dfrac{1}{(n-1)!}\int_0^\infty t^{n-1}e^{-\lambda t}\dfrac{\mathrm{d}}{\mathrm{d}t}(T(t)x)\,\mathrm{d}t\\
&=-\dfrac{1}{(n-2)!}\int_0^\infty t^{n-2}e^{-\lambda t}T(t)x\,\mathrm{d}x\\
&\quad+\dfrac{\lambda}{(n-1)!}\int_0^\infty t^{n-1}e^{-\lambda t}T(t)x\,\mathrm{d}x\\
&=-I_{n-1}(\lambda)x+\lambda I_n(\lambda)x.
\end{aligned}\]For $n=1$, we find instead $I_1(\lambda)Ax=-x+\lambda I_1(\lambda)x$. We can also apply $(T(h)-I)/h$ to \eqref{eq:hille-yosida-resolvent-power-laplace-integral} and go through a similar calculation involving ``differencing by parts.\textquotesingle\textquotesingle\ We find, in the limit $h\to0$, that\[\begin{aligned}
AI_n(\lambda)x&=-I_{n-1}(\lambda)x+\lambda I_m(\lambda)x,\ x\in X,\ n>1,\\
AI_1(\lambda)x&=-x+\lambda I_1(\lambda)x,\ x\in X.
\end{aligned}\]In summary, we find that $I_n(\lambda)(A-\lambda I)=(A-\lambda I)I_n(\lambda)=-I_{n-1}(\lambda)$ for $n>1$ and $I_1(\lambda)(A-\lambda I)=(A-\lambda I)I_1(\lambda)=-I$. This clearly implies that $R_\lambda(A)$ exists and $I_n(\lambda)=(-R_\lambda(A))^n$.
\end{Proof}
\endgroup

We now turn to the proof of sufficiency. For this, we consider the operators
\begin{equation}\label{eq:hille-yosida-implicit-resolvent-power-approximant}
U_n(t)=\left(I-\dfrac{t}{n}A\right)^{-n}.
\end{equation}
It follows from condition~\ref{hyp:hille-yosida-resolvent-powers} that $U_n(t)$ is well defined for $n/t>\omega$, i.e., in particular for sufficiently large $n$. Moreover, the operators $U_n(t)$ are uniformly bounded as $n\to\infty$. For $t>0$, differentiation with respect to $t$ yields
\begin{equation}\label{eq:hille-yosida-approximant-time-derivative}
\dot U_n(t)=A\left(I-\dfrac{t}{n}A\right)^{-n-1}
\end{equation}
in the sense of the operator norm topology. As $t\to0$, we claim that $U_n(t)\to I$ strongly. For this, it suffices to show that $(I-\frac{t}{n}A)^{-1}$ converges to $I$ strongly; see Problem~\ref{ex:strong-operator-products-converge}. Since we know that $(I-\frac{t}{n}A)^{-1}$ is uniformly bounded as $t\to0$, it suffices to show that $(I-\frac{t}{n}A)^{-1}u$ converges for $u$ in a dense set. We now note that, for $u\in\mathcal{D}(A)$,
\begin{equation}\label{eq:hille-yosida-resolvent-initial-strong-limit-bound}
\left\|\left(I-\dfrac{t}{n}A\right)^{-1}u-u\right\|=\left\|\dfrac{t}{n}\left(I-\dfrac{t}{n}A\right)^{-1}Au\right\|\leq Ct\|Au\|,
\end{equation}
which tends to zero as $t\to0$.

We claim that $U_n(t)$ has a strong limit as $n\to\infty$. For that, we estimate $U_n(t)u-U_m(t)u$ for $u$ in the dense set $\mathcal{D}(A^2)=(A-\lambda)^{-1}\mathcal{D}(A)$, $\operatorname{Re}\lambda>\omega$. We note that\[\begin{aligned}
U_n(t)u-U_m(t)u&=\lim_{\epsilon\to0}\int_\epsilon^{t-\epsilon}\dfrac{\mathrm{d}}{\mathrm{d}s}[U_m(t-s)U_n(s)u]\,\mathrm{d}s\\
&=\lim_{\epsilon\to0}\int_\epsilon^{t-\epsilon}[-\dot U_m(t-s)U_n(s)u+U_m(t-s)\dot U_n(s)u]\,\mathrm{d}s.
\end{aligned}\]Using \eqref{eq:hille-yosida-approximant-time-derivative}, one finds, after some algebra, that
\begin{equation}\label{eq:hille-yosida-two-approximants-difference-integral}
\begin{aligned}U_n(t)u-U_m(t)u=\int_0^t&\left(\dfrac{s}{n}-\dfrac{t-s}{m}\right)\left(I-\dfrac{t-s}{m}A\right)^{-m-1}\\&\left(I-\dfrac{s}{n}A\right)^{-n-1}A^2u\,\mathrm{d}s.\end{aligned}
\end{equation}
It follows that
\begin{equation}\label{eq:hille-yosida-approximants-cauchy-bound}
\begin{aligned}\|U_n(t)u-U_m(t)u\|&\leq C\|A^2u\|\int_0^t\dfrac{s}{n}+\dfrac{t-s}{m}\,\mathrm{d}s\\&=\dfrac{Ct^2}{2}\left(\dfrac{1}{n}+\dfrac{1}{m}\right)\|A^2u\|.\end{aligned}
\end{equation}
Hence the strong limit of $U_n(t)$ as $n\to\infty$ exists, and we shall call this limit $T(t)$. Since
\begin{equation}\label{eq:hille-yosida-approximant-exponential-growth-bound}
\|U_n(t)\|\leq M\left(1-\dfrac{t}{n}\omega\right)^{-n},
\end{equation}
we find that $\|T(t)\|\leq M\exp(\omega t)$ in the limit $n\to\infty$.

It remains to be shown that $T(t)$ is indeed a semigroup and that $A$ is its infinitesimal generator. It is easy to see from \eqref{eq:hille-yosida-approximants-cauchy-bound} that the convergence of $U_n(t)u$ to $T(t)u$ is indeed uniform for $t$ in finite intervals. Since $U_n(t)u$ is continuous for every $n$, it follows that $T(t)u$ is continuous for every $u\in\mathcal{D}(A^2)$. Since $\mathcal{D}(A^2)$ is dense and $\|T(t)\|$ is bounded on finite intervals, it follows that $T(t)u$ is continuous for every $u\in X$. Hence $T(t)$ is strongly continuous. It is also obvious that $T(0)=I$, since $U_n(0)=I$ for every $n$. Moreover, we note that the same calculation which led to \eqref{eq:hille-yosida-two-approximants-difference-integral} yields
\begin{equation}\label{eq:hille-yosida-semigroup-composition-derivative-bound}
\begin{aligned}\dfrac{\mathrm{d}}{\mathrm{d}s}[U_n(t-s)U_n(s)u]&=\dfrac{2s-t}{n}\left(I-\dfrac{t-s}{n}A\right)^{-n-1}\\&\quad\left(I-\dfrac{s}{n}A\right)^{-n-1}A^2u\end{aligned}
\end{equation}
for $u\in\mathcal{D}(A^2)$. Letting $n\to\infty$, we find that $T(t-s)T(s)$ is independent of $s$ for $u\in\mathcal{D}(A^2)$ and, by density, for $u\in X$. This is the semigroup property.

\begingroup
\renewenvironment{Proof}{\par}{\par\nobreak\hfill$\blacksquare$\par\vspace{0.1cm}}
\begin{Proof}
By letting $n\to\infty$ in \eqref{eq:hille-yosida-approximant-time-derivative} we find that
\begin{equation}\label{eq:hille-yosida-limit-generator-initial-derivative}
\left.\dfrac{\mathrm{d}}{\mathrm{d}t}T(t)u\right|_{t=0}=Au
\end{equation}
for every $u\in\mathcal{D}(A)$; hence the infinitesimal generator of the semigroup is an extension of $A$. This extension cannot be proper, because according to the growth estimate for $\|T(t)\|$ and the necessity part of this proof, the infinitesimal generator has a resolvent for $\lambda>\omega$. This resolvent must agree with the resolvent of $A$, since a bounded operator defined on the whole space cannot be properly extended.
\end{Proof}
\endgroup

The above proof not only yields the existence of a semigroup; it also yields a practical way of approximating the semigroup. We have
\begin{equation}\label{eq:semigroup-implicit-resolvent-approximation-formula}
T(t)=\lim_{n\to\infty}\left(I-\dfrac{t}{n}A\right)^{-n}.
\end{equation}
This corresponds to the difference scheme
\begin{equation}\label{eq:semigroup-implicit-euler-difference-scheme}
\dfrac{u(t+h)-u(t)}{h}=Au(t+h)
\end{equation}
for solving the equation $\dot u=Au$, which is known as the implicit Euler scheme. Also, it corresponds to formula \eqref{eq:matrix-exponential-resolvent-power-limit} for the matrix exponential. We now feel fully justified in dispensing with separate notations for a semigroup and its infinitesimal generator.

\begin{Definition}
\label{def:semigroup-exponential-notation}
Let $A$ be the infinitesimal generator of a $C_0$-semigroup of bounded linear operators on $X$. Then $\exp(At)$, $t\geq0$ denotes the semigroup generated by $A$.
\end{Definition}

We also introduce the following convenient notation.

\begin{Definition}
\label{def:semigroup-generator-class}
We say that a linear operator $A$ in $X$ is in $G(M,\omega)$ if it satisfies the hypotheses of Theorem~\ref{thm:hille-yosida-generation}.
\end{Definition}

\subsubsection{The Lumer-Phillips Theorem}\label{subsec:lumer-phillips-theorem}

The conditions of the Hille-Yosida theorem are still not easy to verify in applications since they require bounds on all powers of the resolvent. However, the situation is much simpler if $M=1$: If $\|R_\lambda(A)\|\leq(\lambda-\omega)^{-1}$, then obviously $\|R_\lambda(A)^n\|\leq(\lambda-\omega)^{-n}$. Because of the importance of this case, we make the following definition.

\begin{Definition}
\label{def:quasicontraction-contraction-semigroup}
A $C_0$-semigroup $T(t)$ is called a \textbf{quasicontraction semigroup} if $\|T(t)\|\leq\exp(\omega t)$ for some $\omega$. It is called a \textbf{contraction semigroup} if $\|T(t)\|\leq1$.
\end{Definition}

It is obvious that if $A$ generates a quasicontraction semigroup, then $A-\omega I$ generates a contraction semigroup. Actually, every semigroup can be made into a quasicontraction semigroup.

\begin{Theorem}
\label{thm:semigroup-equivalent-quasicontraction-norm}
Let $T(t)$ be a $C_0$-semigroup on $X$ satisfying
\begin{equation}\label{eq:semigroup-equivalent-quasicontraction-norm-growth-hypothesis}
\|T(t)\|\leq M\exp(\omega t).
\end{equation}
Then there is an equivalent norm on $X$ such that, in the operator norm corresponding to this new norm on $X$, we have $\|T(t)\|\leq\exp(\omega t)$.
\end{Theorem}

We omit the proof of this theorem; it can be found, e.g., in \cite{Pa}. In many applications, it is more profitable to seek an appropriate equivalent norm than to try to verify the assumptions of the Hille-Yosida theorem directly. A practical criterion for generators of quasicontraction semigroups is given by the following result, known as the Lumer-Phillips theorem.

\begin{Theorem}{Lumer-Phillips}
\label{thm:lumer-phillips-generation}
Let $H$ be a Hilbert space and let $A$ be a linear operator in $H$ satisfying the following conditions:\begin{enumerate}\item\label{hyp:lumer-phillips-dense-domain} $\mathcal{D}(A)$ is dense.\item\label{hyp:lumer-phillips-quasidissipativity} $\operatorname{Re}(x,Ax)\leq\omega(x,x)$ for every $x\in\mathcal{D}(A)$.\item\label{hyp:lumer-phillips-surjectivity} There exists a $\lambda_0>\omega$ such that $A-\lambda_0I$ is onto.\end{enumerate}Then $A$ generates a quasicontraction semigroup and $\|\exp(At)\|\leq\exp(\omega t)$.
\end{Theorem}

\begingroup
\renewcommand{\endProof}{\par\nobreak\hfill$\blacksquare$\par\vspace{0.1cm}}
\begin{Proof}
We find, for $\lambda>\omega$,
\begin{equation}\label{eq:lumer-phillips-resolvent-lower-bound}
\|(A-\lambda I)x\|\|x\|\geq\operatorname{Re}(x,(\lambda I-A)x)\geq(\lambda-\omega)(x,x),
\end{equation}
i.e., we have $\|(A-\lambda I)x\|\geq(\lambda-\omega)\|x\|$. If we can show that $A-\lambda I$ is onto for every $\lambda>\omega$, then this implies that $R_\lambda(A)$ exists and $\|R_\lambda(A)\|\leq(\lambda-\omega)^{-1}$. In particular, we have $\lambda_0\in\rho(A)$, and hence $A$ is closed. For $\lambda>\omega$, $A-\lambda I$ has a bounded inverse, and hence its range must be closed. But this means that $A-\lambda I$ is semi-Fredholm and its index must be constant on $(\omega,\infty)$.
\end{Proof}
\endgroup

\begin{Example}
\label{ex:self-skew-adjoint-semigroup-generation}
Any self-adjoint operator whose spectrum is bounded from above generates a quasicontraction semigroup. Any skew-adjoint operator generates a contraction semigroup.
\end{Example}

\begin{Example}
\label{ex:absorbing-translation-dissipative-generator}
Let $H=L^2(0,1)$ and let $Au=u^{\prime}$ with domain $\mathcal{D}(A)=\{u\in H^1(0,1)\mid u(1)=0\}$. Clearly $\mathcal{D}(A)$ is dense and it is easy to show that the spectrum of $A$ is empty. Moreover, we have
\begin{equation}\label{eq:absorbing-translation-generator-dissipation-identity}
(u,Au)=\int_0^1u(x)u^{\prime}(x)\,\mathrm{d}x=-\dfrac{1}{2}(u(0))^2\leq0
\end{equation}
for every $u\in\mathcal{D}(A)$; hence $A$ generates a contraction semigroup. Indeed, the semigroup generated by $A$ is the one given in Example~\ref{ex:absorbing-translation-minus-infinite-type}; cf. Problem~\ref{ex:absorbing-translation-generator-characterization}.
\end{Example}

\begin{Definition}
\label{def:quasidissipative-dissipative-quasi-m}
An operator in a Hilbert space which satisfies condition~\ref{hyp:lumer-phillips-quasidissipativity} of Theorem~\ref{thm:lumer-phillips-generation} is called \textbf{quasidissipative}; it is called \textbf{dissipative} if $\omega=0$. An operator which satisfies conditions~\ref{hyp:lumer-phillips-quasidissipativity} and~\ref{hyp:lumer-phillips-surjectivity} is called \textbf{quasi-m-dissipative}.
\end{Definition}

\Needspace{5\baselineskip}\textbf{Problems}

\begin{enumerate}[label=\thesection.\arabic*.,ref=\thesection.\arabic*,start=9,series=book-problems]

\item\label{ex:hille-yosida-heat-dirichlet-existence} Prove the existence of a solution of the heat equation\[u_t=-Lu+f(t),\quad u(0)=u_0,\]where $L$ is the second-order elliptic operator defined in \eqref{eq:second-order-garding-divergence-operator} with Dirichlet boundary conditions. Identify the spaces that the data and solution should occupy.

\item\label{ex:strong-operator-products-converge} Let $A_k\in\mathcal{L}(X,Y)$, $B_k\in\mathcal{L}(Y,Z)$ and assume that $A_k\to A$, $B_k\to B$ as $k\to\infty$ strongly. Show that $B_kA_k\to BA$ strongly.

\item\label{ex:quasicontraction-bounded-operator-perturbation} Let $A\in G(1,\omega)$ and let $B$ be bounded. Show that $A+B\in G(1,\omega+\|B\|)$.

\item\label{ex:translation-opposite-boundary-not-generator} In Example~\ref{ex:absorbing-translation-dissipative-generator}, replace the boundary condition $u(1)=0$ by the condition $u(0)=0$. Show that the resulting operator does \emph{not} generate a $C_0$-semigroup.

\item\label{ex:opposite-generators-inverse-semigroup-group} If both $A$ and $-A$ satisfy the hypotheses of the Hille-Yosida theorem, show that $\exp(At)$ and $\exp((-A)t)$ are inverse to each other. It is then of course natural to write $\exp(-At)$ for $\exp((-A)t)$. Show that the operators $\exp(At)$, $t\in\mathbb{R}$ form a group.

\item\label{ex:skew-adjoint-semigroup-unitary} Let $A$ be skew-adjoint. Show that $\exp(At)$ is unitary.

\end{enumerate}

\subsection{Applications to PDEs}\label{sec:semigroup-pde-applications}

\subsubsection{Symmetric Hyperbolic Systems}\label{subsec:semigroup-symmetric-hyperbolic-systems}

Let $\bm{A}^i(\bm{x})$ be symmetric $p\times p$ matrices defined for $\bm{x}\in\mathbb{R}^m$. We consider the operator
\begin{equation}\label{eq:semigroup-symmetric-hyperbolic-spatial-operator}
\bm{A}\bm{u}(\bm{x})=\bm{A}^i(\bm{x})\dfrac{\partial\bm{u}}{\partial x_i}
\end{equation}
in the function space $L^2(\mathbb{R}^m)$. We assume that $\bm{A}^i\in C_b^1(\mathbb{R}^m)$. It is no simple matter to characterize the domain for $\bm{A}$. Instead, we take a copout. We first define $\bm{A}$ on $H^1(\mathbb{R}^m)$ and then take the closure. We claim

\begin{Theorem}
\label{thm:symmetric-hyperbolic-semigroup-generation}
$\bm{A}$ is the infinitesimal generator of a $C_0$-semigroup.
\end{Theorem}

\begingroup
\renewcommand{\endProof}{\par}
\begin{Proof}
We shall use the Lumer-Phillips theorem. Since the domain of $\bm{A}$ includes $H^1(\mathbb{R}^m)$, it is clear that $\bm{A}$ is densely defined. Moreover, for $\bm{u}\in H^1$, an integration by parts yields\[\begin{aligned}
(\bm{u},\bm{A}\bm{u})&=\int_{\mathbb{R}^m}\bm{u}(\bm{x})\cdot\bm{A}^i(\bm{x})\dfrac{\partial\bm{u}}{\partial x_i}\,\mathrm{d}\bm{x}\\
&=-\dfrac{1}{2}\int_{\mathbb{R}^m}\bm{u}\cdot\dfrac{\partial\bm{A}^i}{\partial x_i}\bm{u}\,\mathrm{d}\bm{x}.
\end{aligned}\]Hence the second condition of the Lumer-Phillips theorem holds for every $\bm{u}\in H^1(\mathbb{R}^m)$, and we can take limits to show it holds for every $\bm{u}\in\mathcal{D}(\bm{A})$.
\end{Proof}
\endgroup

The main part of the work is in verifying the third condition, namely, that $\bm{A}-\lambda I$ is onto for sufficiently large $\lambda\in\mathbb{R}$. From the proof of the Lumer-Phillips theorem it is clear that it suffices to show that the range of $\bm{A}-\lambda I$ is dense. For this purpose, we first consider a new operator $\widetilde{\bm{A}}$, which agrees with $\bm{A}$ except for the domain. Namely, we take the domain of $\widetilde{\bm{A}}$ to be the set of all $\bm{u}\in L^2$ such that $\bm{A}^i\partial\bm{u}/\partial x_i$ (interpreted as an element of $H^{-1}(\mathbb{R}^m)$) is also in $L^2(\mathbb{R}^m)$. It is clear that $\widetilde{\bm{A}}$ is an extension of $\bm{A}$. We next use a Galerkin argument to show that $\widetilde{\bm{A}}-\lambda I$ is onto for $\lambda$ sufficiently large. For this, let $\{\bm{\phi}_n\}$, $n\in\mathbb{N}$, be a complete orthogonal system in $H^1(\mathbb{R}^m)$. Let $X_n$ be the span of $\{\bm{\phi}_1,\ldots,\bm{\phi}_n\}$, and let $P_n$ be the orthogonal projection from $L^2(\mathbb{R}^m)$ onto $X_n$. We now consider the equation
\begin{equation}\label{eq:symmetric-hyperbolic-resolvent-galerkin-equation}
(\bm{\phi},(\bm{A}-\lambda I)\bm{u}_n)=(\bm{\phi},\bm{f})\quad\forall\bm{\phi}\in X_n,
\end{equation}
where we seek a solution $\bm{u}_n\in X_n$. Using the fact that $\bm{A}$ is quasidissipative, it is easy to see that such a solution $\bm{u}_n$ exists and that the $L^2$-norm of $\bm{u}_n$ is bounded independently of $n$. Hence a subsequence of the $\bm{u}_n$ converges weakly, and the limit is easily shown to be a solution of $(\widetilde{\bm{A}}-\lambda I)\bm{u}=\bm{f}$.\par It remains to be shown that $\widetilde{\bm{A}}$ is indeed $\bm{A}$. That is, given $\bm{u},\bm{f}\in L^2(\mathbb{R}^m)$ such that $\widetilde{\bm{A}}\bm{u}=\bm{f}$, we have to find $\bm{u}_n\in H^1(\mathbb{R}^m)$, $\bm{f}_n\in L^2(\mathbb{R}^m)$ such that $\bm{A}\bm{u}_n=\bm{f}_n$ and $\bm{u}_n\to\bm{u}$, $\bm{f}_n\to\bm{f}$ in $L^2(\mathbb{R}^m)$. For this, we shall need the following result, known as Friedrichs\textquotesingle\ lemma.

\begin{Lemma}{Friedrichs}
\label{lem:friedrichs-commutator-approximation}
For $\epsilon>0$, there exists a mapping $Q_\epsilon$ from $L^2(\mathbb{R}^m)$ into $H^1(\mathbb{R}^m)$ such that the following properties hold:\begin{enumerate}\item For every $\bm{f}\in L^2(\mathbb{R}^m)$, $Q_\epsilon\bm{f}\to\bm{f}$ in $L^2(\mathbb{R}^m)$ as $\epsilon\to0$.\item The operator $\bm{A}Q_\epsilon-Q_\epsilon\bm{A}$ can be extended to a bounded operator in $L^2(\mathbb{R}^m)$. Moreover, for any fixed $\bm{u}\in\mathcal{D}(\widetilde{\bm{A}})$, $\bm{A}Q_\epsilon\bm{u}-Q_\epsilon\widetilde{\bm{A}}\bm{u}\to0$ in $L^2(\mathbb{R}^m)$.\end{enumerate}
\end{Lemma}

With this lemma, the completion of the proof is immediate. We simply set $\bm{u}_\epsilon=Q_\epsilon\bm{u}$ and $\bm{f}_\epsilon=\bm{A}\bm{u}_\epsilon$. Then
\begin{equation}\label{eq:symmetric-hyperbolic-friedrichs-graph-approximation}
\bm{f}_\epsilon=Q_\epsilon\widetilde{\bm{A}}\bm{u}+(\bm{A}Q_\epsilon-Q_\epsilon\widetilde{\bm{A}})\bm{u}.
\end{equation}
As $\epsilon\to0$, the first term on the right converges to $\widetilde{\bm{A}}\bm{u}=\bm{f}$, and the second term converges to zero.

\begin{Proof}
It thus remains to prove Friedrichs\textquotesingle\ lemma. Let $k(\bm{x})$ be a non-negative test function supported on the ball $|\bm{x}|\leq1$ and having unit integral. We define $k_\epsilon(\bm{x})=\epsilon^{-m}k(\bm{x}/\epsilon)$, and we define $Q_\epsilon$ to be convolution with $k_\epsilon$. Then, for every $\epsilon>0$, $Q_\epsilon$ takes $L^2(\mathbb{R}^m)$ to $H^1(\mathbb{R}^m)$ (in fact to $C^\infty(\mathbb{R}^m)$); moreover, Lemma~\ref{lem:integral-kernel-row-column-boundedness} implies that $Q_\epsilon$ has norm less than or equal to 1 as an operator from $L^2(\mathbb{R}^m)$ to itself. It is easy to show that $Q_\epsilon\bm{u}\to\bm{u}$ for a dense subset of $L^2(\mathbb{R}^m)$ (e.g., for $\bm{u}\in\mathcal{D}(\mathbb{R}^m)$), and together with the uniform boundedness of $Q_\epsilon$, we conclude that $Q_\epsilon$ converges to the identity strongly, i.e., property 1 of the lemma holds.\par To verify property 2, one calculates
\begin{equation}\label{eq:friedrichs-mollifier-commutator-integral-identity}
\begin{aligned}
\bm{A}Q_\epsilon\bm{u}-Q_\epsilon\widetilde{\bm{A}}\bm{u}&=\int_{\mathbb{R}^m}k_\epsilon(\bm{x}-\bm{y})(\bm{A}^i(\bm{x})-\bm{A}^i(\bm{y}))\dfrac{\partial\bm{u}}{\partial y_i}(\bm{y})\,\mathrm{d}\bm{y}\\
&=\int_{\mathbb{R}^m}k_\epsilon(\bm{x}-\bm{y})\dfrac{\partial\bm{A}^i}{\partial y_i}(\bm{y})\bm{u}(\bm{y})\,\mathrm{d}\bm{y}\\
&\quad+\int_{\mathbb{R}^m}\dfrac{\partial}{\partial x_i}(k_\epsilon(\bm{x}-\bm{y}))(\bm{A}^i(\bm{x})-\bm{A}^i(\bm{y}))\bm{u}(\bm{y})\,\mathrm{d}\bm{y}.
\end{aligned}
\end{equation}
As $\epsilon\to0$, the first expression on the right-hand side clearly tends to zero for smooth enough $\bm{u}$ (e.g., $\bm{u}\in\mathcal{D}(\mathbb{R}^m)$). Hence it suffices to show that $\bm{A}Q_\epsilon-Q_\epsilon\widetilde{\bm{A}}$ is uniformly bounded for $\epsilon\to0$ as an operator from $L^2(\mathbb{R}^m)$ to itself. For this purpose, we use the second expression on the right-hand side of \eqref{eq:friedrichs-mollifier-commutator-integral-identity}. We can see from the proof of Lemma~\ref{lem:integral-kernel-row-column-boundedness} that the first term in this expression has an operator norm bounded by
\begin{equation}\label{eq:friedrichs-commutator-coefficient-derivative-norm-bound}
\sup_{\bm{y}\in\mathbb{R}^m}\left|\dfrac{\partial\bm{A}^i}{\partial y_i}(\bm{y})\right|\times\int_{\mathbb{R}^m}|k_\epsilon(\bm{x})|\,\mathrm{d}\bm{x},
\end{equation}
and in this expression the first factor is independent of $\epsilon$ and the second is 1. The operator norm for the second term can, again by Lemma~\ref{lem:integral-kernel-row-column-boundedness}, be estimated by
\begin{equation}\label{eq:friedrichs-commutator-kernel-derivative-norm-bound}
\sup_{\bm{x}\in\mathbb{R}^m}\int_{\mathbb{R}^m}\left|\dfrac{\partial}{\partial x_i}k_\epsilon(\bm{x}-\bm{y})\right|\left|\bm{A}^i(\bm{x})-\bm{A}^i(\bm{y})\right|\,\mathrm{d}\bm{y}.
\end{equation}
Since the support of $k_\epsilon$ is contained in $|\bm{x}|\leq\epsilon$, we can estimate $|\bm{A}^i(\bm{x})-\bm{A}^i(\bm{y})|$ by a constant times $\epsilon$. On the other hand, the $L^1$-norm of $\partial k_\epsilon/\partial x_i$ is a constant times $1/\epsilon$. Hence \eqref{eq:friedrichs-commutator-kernel-derivative-norm-bound} is bounded by a constant independent of $\epsilon$.
\end{Proof}

\subsubsection{The Wave Equation}\label{subsec:semigroup-wave-equation}

We consider the second-order equation
\begin{equation}\label{eq:semigroup-wave-variable-coefficient-equation}
u_{tt}=\dfrac{\partial}{\partial x_i}\left(a_{ij}(\bm{x})\dfrac{\partial u}{\partial x_j}\right),\ \bm{x}\in\Omega,\ t>0,
\end{equation}
with boundary condition
\begin{equation}\label{eq:semigroup-wave-dirichlet-boundary}
u(\bm{x},t)=0,\ \bm{x}\in\partial\Omega,
\end{equation}
and initial conditions
\begin{equation}\label{eq:semigroup-wave-displacement-velocity-initial-data}
u(\bm{x},0)=u_0(\bm{x}),\ u_t(\bm{x},0)=u_1(\bm{x}).
\end{equation}
Here we assume that $\Omega$ is a bounded domain in $\mathbb{R}^m$ with smooth boundary and that the matrix $a_{ij}$ is symmetric, strictly positive definite and of class $C^1(\overline\Omega)$.

We denote by $A$ the operator represented by the right-hand side of \eqref{eq:semigroup-wave-variable-coefficient-equation}. We view $A$ as an operator in $L^2(\Omega)$ with domain $H^2(\Omega)\cap H_0^1(\Omega)$. In order to apply semigroup theory, we must first transform the second-order equation \eqref{eq:semigroup-wave-variable-coefficient-equation} to a first-order system. For this purpose, we set $v=u_t$, and we rewrite \eqref{eq:semigroup-wave-variable-coefficient-equation} as
\begin{equation}\label{eq:semigroup-wave-first-order-operator-matrix}
\begin{pmatrix}\dot u\\\dot v\end{pmatrix}=\mathcal{A}\begin{pmatrix}u\\v\end{pmatrix}:=\begin{pmatrix}0&I\\A&0\end{pmatrix}\begin{pmatrix}u\\v\end{pmatrix}.
\end{equation}
We regard $\mathcal{A}$ as an operator in the space $X=H_0^1(\Omega)\times L^2(\Omega)$ with domain $\mathcal{D}(\mathcal{A})=(H_0^1(\Omega)\cap H^2(\Omega))\times H_0^1(\Omega)$. We claim

\begin{Theorem}
\label{thm:semigroup-wave-generation}
The operator $\mathcal{A}$ generates a $C_0$-semigroup.
\end{Theorem}

\begin{Proof}
We again use the Lumer-Phillips theorem. For this purpose, we consider the following inner product on $X$, which is equivalent to the usual one:
\begin{equation}\label{eq:semigroup-wave-energy-inner-product}
\left[\begin{pmatrix}u\\v\end{pmatrix},\begin{pmatrix}f\\g\end{pmatrix}\right]=\int_\Omega a_{ij}(\bm{x})\dfrac{\partial u}{\partial x_i}\dfrac{\partial f}{\partial x_j}+vg\,\mathrm{d}\bm{x}.
\end{equation}
It is easy to see that
\begin{equation}\label{eq:semigroup-wave-generator-energy-skew-pairing}
\left[\begin{pmatrix}u\\v\end{pmatrix},\mathcal{A}\begin{pmatrix}u\\v\end{pmatrix}\right]=0
\end{equation}
for $(u,v)\in\mathcal{D}(\mathcal{A})$. Moreover, it is clear that $\mathcal{D}(\mathcal{A})$ is dense. We claim that every positive $\lambda$ is in the resolvent set of $\mathcal{A}$. For this, we need to consider the problem
\begin{equation}\label{eq:semigroup-wave-resolvent-coupled-equations}
v-\lambda u=f,\ Au-\lambda v=g,\ f\in H_0^1(\Omega),\ g\in L^2(\Omega).
\end{equation}
These two equations can be combined into the one equation
\begin{equation}\label{eq:semigroup-wave-resolvent-reduced-elliptic-equation}
Au-\lambda^2u=g+\lambda f.
\end{equation}
The solvability of this equation follows from our results on elliptic boundary-value problems in Chapter~\ref{chap:linear-elliptic-equations}.
\end{Proof}

\subsubsection{The Schrödinger Equation}\label{subsec:semigroup-schrodinger-equation}

The Schrödinger equation is the fundamental equation of quantum mechanics. It reads
\begin{equation}\label{eq:semigroup-schrodinger-real-potential-equation}
u_t=i(\Delta u-V(\bm{x})u),\ \bm{x}\in\mathbb{R}^m.
\end{equation}
Here $i$ is the imaginary unit, i.e., $u$ is complex-valued. However, the function $V(\bm{x})$ is real-valued. In a formal sense, the operator on the right-hand side of \eqref{eq:semigroup-schrodinger-real-potential-equation} is skew-symmetric, and we expect it to be skew-adjoint if the domain is chosen correctly. In this case, it follows immediately that the operator generates a $C_0$-semigroup.\par A major part of the effort in mathematical quantum mechanics goes into proving that this operator is skew-adjoint. How difficult this is depends, obviously, on how badly $V$ is allowed to behave. The case where $V$ is bounded is basically trivial, but of little physical interest. Somewhat more interesting are potentials with point singularities. We prove here the following result.

\begin{Theorem}
\label{thm:schrodinger-ltwo-potential-skew-adjointness}
Assume that $m\leq3$ and $V\in L^2(\mathbb{R}^m)$. Then the operator $A$ given by $Au=i(\Delta u-V(\bm{x})u)$ with domain $\mathcal{D}(A)=H^2(\mathbb{R}^m)$ is skew-adjoint.
\end{Theorem}

\begin{Proof}
We note that
\begin{equation}\label{eq:schrodinger-ltwo-potential-multiplication-bound}
\|Vu\|_2\leq\|V\|_2\|u\|_\infty\leq C\|V\|_2\|u\|_{m/2+\delta,2}\leq C\|V\|_2\|u\|_{2,2}.
\end{equation}
Here $\delta$ is any sufficiently small positive number. This shows that $A$ is indeed defined on $H^2(\mathbb{R}^m)$. Moreover, a straightforward calculation involving Fourier transforms shows that we can improve on \eqref{eq:schrodinger-ltwo-potential-multiplication-bound} as follows. For every $\epsilon>0$, there is a constant $C(\epsilon)$ such that
\begin{equation}\label{eq:schrodinger-potential-small-relative-interpolation-bound}
\|u\|_{m/2+\delta,2}\leq\epsilon\|u\|_{2,2}+C(\epsilon)\|u\|_2.
\end{equation}
We can use this to show that $A$ is closed; see Problem~\ref{ex:schrodinger-potential-operator-closedness}. It is clear that $A$ is skew-symmetric. To show it is skew-adjoint, we have to prove that $A\pm I$ is surjective. This follows by considering the family of operators $i(\Delta-tV)\pm I$ for $t\in[0,1]$ and using the homotopy invariance of the Fredholm index.
\end{Proof}

Clearly, we can more generally consider a $V$ which is the sum of an $L^2$-function and a bounded function. This includes, for example, a Coulomb potential ($V=-1/r$) in $\mathbb{R}^3$.

\Needspace{5\baselineskip}\textbf{Problems}

\begin{enumerate}[label=\thesection.\arabic*.,ref=\thesection.\arabic*,start=15,series=book-problems]

\item\label{ex:semigroup-transport-and-wave-existence-statements} State an existence theorem for the equation $u_t=a(x,y)u_x+b(x,y)u_y+c(x,y)u+f(x,y,t)$, $(x,y)\in\mathbb{R}^2$, $t>0$, with initial condition $u(x,y,0)=u_0(x,y)$. Similarly, state an existence theorem for the wave equation.

\item\label{ex:strict-hyperbolic-system-symmetric-transformation} Consider the problem $\bm{u}_t=\bm{A}(x)\bm{u}_x$, where the matrix $\bm{A}$ has real and distinct eigenvalues. Show that this problem can always be transformed to one with a symmetric matrix to which the results of Section~\ref{subsec:semigroup-symmetric-hyperbolic-systems} are applicable. What difficulty arises in several space dimensions?

\item\label{ex:second-order-drift-operator-hthree-domain-generation} In $L^2(\mathbb{R})$, let $Au=u_{xx}+a(x)u_x$ with domain $\mathcal{D}(A)=H^3(\mathbb{R})$. Assume that $a\in C_b^1(\mathbb{R})$. Prove that $A$ generates a $C_0$-semigroup.

\item\label{ex:wave-neumann-source-unresolved-refwhyp} Discuss the wave equation with Neumann conditions in a manner similar to Section \texttt{refWHYP}.

\item\label{ex:schrodinger-potential-operator-closedness} Use \eqref{eq:schrodinger-potential-small-relative-interpolation-bound} to show that the operator $A$ in Theorem~\ref{thm:schrodinger-ltwo-potential-skew-adjointness} is closed.

\end{enumerate}
